% 中文请用ctexbeamer, 个人觉得行距linespread宽一点看起来比较舒服。XeLaTeX编译。
\documentclass[aspectratio=1610,linespread=1.4,t]{ctexbeamer}

% If this is an English-only deck, you don't need ctexbeamer nor the wider linespread
% \documentclass[aspectratio=1610,t]{beamer}
% But load fontspec so that you can load OTF/TTF fonts if you want to!
% \usepackage{fontspec} 

% beamerthemeMirage.sty v1.0 2024/04/14
% by LianTze Lim (liantze@gmail.com)
% A beamer theme inspired by a musical album art
% 基于《蜃楼》歌曲海报色调的beamer主题
\usetheme{Mirage}           % Default = dark mode
% \usetheme[light]{Mirage}  % Alternative light mode 

% The glowy light in the footnline
% \setlength{\MirageGlowRadius}{0.25ex}

% Mirage theme loads the fontawesome7 package

\title{Mirage Theme 蜃楼}
\subtitle{此主题为基于《蜃楼》歌曲海报的二创}
\author{林莲枝}
\institute{（可选填）某某大学某某单位}
\date{2026/08/27}

% These fonts will work with pdflatex, xelatex, lualatex. If you're using ctexbeamer it's best to compile in xelatex.
\usepackage[lining,tabular]{carlito}
\usepackage{caladea}

% But unicode-math and TTF/OTF fonts will only work with xelatex, lualatex
\usepackage{unicode-math}
% \setmathfont{STIX Two Math}
% \setmathfont{Erewhon Math}
\setmathfont{Fira Math}

\begin{document}

\frame{\maketitle}

\section{简介}
\begin{frame}
\frametitle{一张幻灯片示例}

作为示例的一些文字。

\begin{itemize}
    \item  一些文字，日月盈昃，辰宿列张。
    \begin{itemize}
        \item 日月盈昃，辰宿列张，此为第二级项目符号列表。
        \item 日月盈昃，辰宿列张，此为第二级项目符号列表。
    \end{itemize}
    \item 一些文字，日月盈昃，辰宿列张。
\end{itemize}

\end{frame}

\begin{frame}
\frametitle{再来一张}
\begin{enumerate}
    \item  一些文字，日月盈昃，辰宿列张。
    \begin{itemize}
        \item 日月盈昃，辰宿列张，此为第二级有序列表。
        \item 日月盈昃，辰宿列张，此为第二级有序列表。
    \end{itemize}
\item 再来一句\alert{刻骨铭心的口号}。
\end{enumerate}
\end{frame}

\section{其它示范}

\begin{frame}[c]
\frametitle{拉栏引语}
\begin{columns}
\begin{column}{.4\textwidth}
\begin{pullquote}
    Can it be real\\
    The world is a mirage
\end{pullquote}
\end{column}

\begin{column}{.53\textwidth}
\setbeamercolor{pullquote}{fg=MirageBlue}
\renewcommand{\MiragePullquoteOpen}{\hskip-.2\ccwd『}
\begin{pullquote}
在电幻的荒丘　寻真实的绿洲\\
渺小得如蜉蝣　也仰望着宇宙
\end{pullquote}
\end{column}
\end{columns}
\end{frame}


% \renewcommand{\MirageFrametitlePrefix}{\faDove}
% \renewcommand{\MirageProofPrefix}{\faSquareCheck[regular]}
% \renewcommand{\MirageTheoremPrefix}{\faGears}

\begin{frame}[allowframebreaks]{各种block}
    
    \begin{exampleblock}{算了我也不知道在写什么，do you?}
    Now solve $x = \frac{-b \pm \sqrt{b^2 -4ac}}{2a}$. 对各位同学来说应该挑战不大。
    \end{exampleblock}

\begin{columns}[T]
\column{0.48\textwidth}
    \begin{alertblock}{算了我也不知道在写什么，do you?}
    \[ x = \frac{-b \pm \sqrt{b^2 -4ac}}{2a}, \quad\therefore \alpha \neq \Omega \]
    \end{alertblock}
    
\column{0.48\textwidth}
    \begin{block}{算了我也不知道在写什么，do you?}
    \[ x = \frac{-b \pm \sqrt{b^2 -4ac}}{2a}, \quad\therefore \alpha \neq \Omega \]
    \end{block}
\end{columns}
    
    \begin{proof}
    显而易见，$1+1=2$.
    \end{proof}

    \begin{theorem}
    有一件很美好的事情将要发生，它终会发生。
    \end{theorem}
    
    \begin{definition}
    有一件很美好的事情将要发生，它终会发生。
    \end{definition}
\end{frame}


\end{document}
